\documentclass[12pt]{article}
\usepackage[utf8]{inputenc}

\usepackage{mystyle}

\title{Nullstellensatz}
\author{Joseph Sullivan}
\date{July 2022}

\begin{document}

\maketitle

\section{Introduction}

These notes are a paraphrase/clarification/obfuscation of Gathmann's notes on Commutative Algebra, covering most of the essentials to proving Hilbert's Nullstellensatz.

\begin{defn}
    Let $K$ be a field. We define
    \begin{itemize}
        \item $\A_K^n = \A^n = K^n$.
        \item $I(a) = (x_1-a_1,\dots,x_n-a_n) \teq K[x_1,\dots,x_n]$,
        \item For $S\subseteq K[x_1,\dots,x_n]$, $V(S) = \{a\in \A^n \mid f(a)=0 \text{ for all } f\in S\}$
        \item For $X\subseteq K^n$, $I(X) = \{f\in K[x_1,\dots,x_n] \mid f(a)=0 \text{ for all } a\in X\}$
    \end{itemize}
    In words, $V(S)$ is the set of points where all of $S$ is zero, and $I(X)$ are the functions that are zero on at least $X$.
\end{defn}

\section{Cayley-Hamilton}

\begin{thm}[Cayley-Hamilton]\label{cayham}
    Let $M$ be a finitely generated $R$-module, $I$ an ideal of $R$, and $\phi: M\to M$ an $R$-module homomorphism with $\phi(M)\subseteq IM$. Then there is a monic polynomial
    \[\chi = x^n + a_{n-1} x^{n-1} + \cdots + a_0 \in R[x]\]
    with $a_0,\dots,a_{n-1}\in I$ and
    \[\chi(\phi) \defeq \phi^n + a_{n-1}\phi^{n-1} + \cdots + a_0 \id = 0 \quad \in \Hom_R(M,M)\]
    where $\phi^i$ denotes the $i$-fold composition of $\phi$ with itself.
\end{thm}

\begin{proof}
    Let $m_1,\dots,m_n$ be generators of $M$. We can then represent $\phi$ in this generating set, i.e., write
    \[\phi(m_i) = \sum_{j=1}^n a_{i, j} m_i\]
    for some $a_{i,j} \in I$, so we can write
    \[
        \begin{bmatrix}
            \phi(m_1)\\ \vdots \\ \phi(m_n)
        \end{bmatrix} =
        \begin{bmatrix}
            a_{1,1} & \cdots & a_{1,n}\\
            \vdots & \ddots & \vdots\\
            a_{n,1} & \cdots & a_{n,n}
        \end{bmatrix}
        \begin{bmatrix}
            m_1\\ \vdots \\ m_n
        \end{bmatrix}.
    \]
    We denote this matrix by $A \in M_n(R)$, and we denote our column of $m_i$ by $\mathbf{m} \in M^n$. Upgrading our module $M$ to an $R[x]$-module by defining $x\cdot m = \phi(m)$, we also get
    \[
        \begin{bmatrix}
            \phi(m_1)\\ \vdots \\ \phi(m_n)
        \end{bmatrix} =
        \begin{bmatrix}
            x &  & 0 \\
             & \ddots & \\
            0 &  & x
        \end{bmatrix}
        \begin{bmatrix}
            m_1\\ \vdots \\ m_n
        \end{bmatrix} = xI_n \mathbf{m}
    \]
    Thus, we have
    \[(xI_n - A)\mathbf{m} = 0.\]
    Multiplying by the adjugate matrix in $M_n(R[x])$, we get
    \[\det(xI_n - A)\mathbf{m} = 0.\]
    Expanding the determinant, we get that $\det(xI_n - A)$ is a monic polynomial in $R[x]$ with non-leading coefficients in $\<a_{i,j}\> \subseteq I$. Then, since it annihilates $m_1,\dots,m_n$ which are generators for $M$, so it annihilates all of $M$. Thus, $\chi := \det(xI_n - A)$ is the zero homomorphism in $\Hom_R(M,M)$.
\end{proof}

\begin{rmk}
    If we take $R=K$ a field and $I=K$, $M=V$ a finite dimensional $K$ vector space, and $m_i=e_i$ a $K$-basis for $V$---then we get the classical Cayley-Hamilton with $\chi = p_\phi(\lambda) = \det(\lambda I - \phi)$.
\end{rmk}

\section{Localization}

\begin{prop}\label{getpoint}
    Let $S$ be a multiplicative subset in a ring $R$, and let $I\teq R$ be an ideal with $I\cap S = \emptyset$. Then, there exists a prime ideal $P$ of $R$ containing $I$ with $P\cap S = \emptyset$ and $S^{-1} P$ is a maximal ideal in $S^{-1} R$.
\end{prop}

When we take $S=\{1,f,f^2,\dots\}$, we can interpret this theorem geometrically in scheme language. It tells us in particular that $D(f)\cap V(I)$ is nonempty.

\begin{proof}
    The ideal $S^{-1} I \teq S^{-1} R$ is proper because $I\cap S = \emptyset$. In more detail, otherwise we would have $1\in S^{-1} I$, so there exists $a\in I, s\in S$ so that $\frac{a}{s}=1$. Then, this implies there exists $v\in S$ such that $va=vs$. However, this element would be in $I\cap S$.
    
    Then, since $S^{-1} I$ is proper, it can be extended to a maximal ideal in $S^{-1} R$. The contraction $P$ will then be a prime (not necessarily maximal) ideal disjoint from $S$.
\end{proof}

\section{Integrality}

\begin{defn}[Integral and finite ring extensions]\;
    \begin{itemize}
        \item If $R\subseteq R'$, we call $R'$ an \textbf{extension ring} of $R$, and we call $R\subseteq R'$ a ring extension.
        \item An element $a\in R'\supseteq R$ is called \textbf{integral} over $R$ if there is a monic polynomial $f\in R[x]$ with $f(a)=0$. We say $R'$ is \textbf{integral} over $R$ if every elements of $R'$ is integral over $R$.
        \item A ring extension $R\subseteq R'$ is called \textbf{finite} if $R'$ is finitely generated as an $R$-module.
    \end{itemize}
\end{defn}

\begin{rmk}
    For integrality, we require the existence of a \textit{monic} polynomial because for purposes of finite generation we want to be able to write $a^n$ in terms of lower powers of $a$ for sufficiently large $n$. In a field, we can always make polynomials monic to accomplish this, but for rings we need the monic condition.
\end{rmk}

\begin{prop}\label{integral}
    An extension ring $R'$ is finite over $R$ if and only if $R'=R[a_1,\dots,a_n]$ for integral elements $a_1,\dots, a_n\in R'$ over $R$.
    
    Moreover, in this case the whole ring extension $R\subseteq R'$ is integral.
\end{prop}

\begin{proof}
    Assume $R' = \<a_1,\dots,a_n\>$ (is finitely generated) as an $R$-module. In particular, $R'=R[a_1,\dots,a_n]$. We now just need to show that $R'$ is integral over $R$ to get the ``moreover'' and that $a_1,\dots,a_n$ are integral over $R$. Let $a\in R'$. Applying Cayley-Hamltion theorem \ref{cayham} to the $R$-module homomorphism
    \begin{align*}
        m_a: R' &\longrightarrow R'\\
        x &\longmapsto ax
    \end{align*}
    to get a monic polynomial equation
    \[m_a^k + c_{k-1} m_a^{k-1} + \cdots + c_0 = 0,\]
    for some $c_0,\dots,c_{k-1}\in R$, which applied to $1\in R$ gives us
    \[a^k + c_{k-1} a^{k-1} + \cdots + c_0 = 0.\]
    Thus, $a$ is integral over $R$.
    
    Now, assume $R' = R[a_1,\dots,a_n]$ with $a_1,\dots,a_n$ integral over $R$. Thus, each $a_i$ satisfies a monic relation of degree $d_i$. Therefore, we can write every element of $R'$ as a sum
    \[\sum_{0\leq k_i \leq d_i} c_{k_1,\dots,k_n} a_1^{k_1}\cdots a_n^{k_n}\]
    since the homomorphism $R[x_1,\dots,x_n]\to R[a_1,\dots,a_n]$ is surjective and we can rewrite any higher degree terms in terms of lower degrees using our monic relation. Thus, $R'=\<a_1^{k_1}\cdots a_n^{k_n}\>$ for $k_i=0,\dots,n$ as an $R$-module. Therefore, $R'$ is a finitely generated module over $R$.
\end{proof}

\begin{lem}\label{transint}
    Let $R\subseteq R' \subseteq R''$ be ring extensions.
    \begin{enumerate}[(a)]
        \item If $R\subseteq R'$ and $R'\subseteq R''$ are finite, then so is $R\subseteq R''$.
        \item If $R\subseteq R'$ and $R'\subseteq R''$ are integral, then so is $R\subseteq R''$.
    \end{enumerate}
\end{lem}

\begin{proof}[Proof of (a).]
    By proposition \ref{integral}, we can write $R'=R[a_1,\dots,a_n]$ and $R''=R'[b_1,\dots,b_m]$. Then, we can write $R''=R[a_ib_j \mid i=1,\dots,n, j=1,\dots,m]$, so $R''$ is finite over $R$.
\end{proof}

\begin{proof}[Proof of (b).]
    Let $a\in R''$. Since $R''$ is integral over $R'$, there exists monic polynomial
    \[a^n + c_{n-1} a^{n-1} + \cdots + c_1 a^1 + c_0 = 0,\]
    for $c_0,\dots,c_{n-1} \in R'$, so
    \[a\in R[c_0,\dots,c_{n-1}].\]
    Since $R'$ is integral over $R$, each of $c_0,\dots,c_{n-1}$ are integral over $R$. Therefore, $R[c_0,\dots,c_{n-1}]$ is a finite extension of $R$, so $a\in R[c_0,\dots,c_{n-1}]$ is integral over $R$. Thus, $R''$ is integral over $R$.
\end{proof}

\begin{lem}\label{intquot}
    Let $R\subseteq R'$ be an integral ring extension. If $I\teq R'$, then $R'/I$ is an integral ring extension of $R/I$.
\end{lem}

\begin{proof}
    Let $\overline{a}\in R'$. Then, $a\in R'$ satisfies some monic relation
    \[a^n + c_{n-1}a^{n-1} + \cdots + c_1 a + c_0 = 0\]
    for $c_0,\dots,c_{n-1}\in R$, so $\overline{a}$ satisfies the relation
    \[\overline{a}^n + \overline{c_{n-1}}\overline{a}^{n-1} + \cdots + \overline{c_1} \overline{a} + \overline{c_0} = 0\]
    for $\overline{c_0},\dots,\overline{c_{n-1}} \in R/(I\cap R)$ (identified with subring of $R'/I$), so $\overline{a}$ is integral over $R/(I\cap R)$.
\end{proof}

\begin{prop}[Lying Over]\label{lyingover}
    Let $R\subseteq R'$ be a ring extension, and $P\teq R$ prime.
    \begin{enumerate}[(a)]
        \item There is a prime ideal $P'\teq R'$ with $P'\cap R = P$ if and only if $PR' \cap R \subseteq P$.
        \item If $R\subseteq R'$ is integral, then this is always the case.
    \end{enumerate}
    We say in this case that $P'$ is lying over $P$.
\end{prop}

The condition in (a) is saying that the extension $P^e$ isn't too large, i.e., that when we contract to get $P^{ec}$ we are bounded by $P$. Geometrically, (b) says that the map $\Spec R' \to \Spec R$ is surjective when $R\subseteq R'$ is integral.

\begin{proof}[Proof of (a).]
    Assume there exists $P' \teq R'$ such that $P'\cap R = P$. Then,
    \[PR' \cap R = (P'\cap R)R' \cap R \subseteq P'R' \cap RR' \cap R = P' \cap R \cap R = P\]
    (watch out that product and intersection need not commute exactly).
    
    Now, assume that $PR'\cap R \subseteq P$. We are looking for a prime ideal $P'\teq R$ containing $PR'$ and not much more. By proposition \ref{getpoint}, since for $S=R\setminus P$, we have
    \[PR' \cap S = PR'\cap R\setminus P = \emptyset,\]
    there exists a prime ideal $P'\teq R'$ containing $PR'\supseteq P$ with $P' \cap S = \emptyset$. This means $P'\cap (R\setminus P) = \emptyset$, and since $P\subseteq P'$, this means $P'\cap R = P$.
\end{proof}

\begin{proof}[Proof of (b).]
    We show that $PR' \cap R \subseteq P$. Let $a\in PR' \cap R$. In particular, $a\in PR'$, so we have $R$-module homomoprhism
    \begin{align*}
        m_a: R' &\longrightarrow PR'\\
        x &\longmapsto ax.
    \end{align*}
    By Cayley Hamilton theorem \ref{cayham}, there exist $c_0,\dots,c_{n-1} \in P$ such that
    \begin{align*}
        m_a^n + c_{n-1}m_a^{n-1} + \cdots + c_1m_a + c_0 &= 0\\
        a^n = -(c_{n-1}a^{n-1} + \cdots + c_1 a + c_0).
    \end{align*}
    Since $a\in R$, the left hand expression is in $RP=P$, so $a^n\in P$. Since $P$ is prime, $a^n\in P$ implies $a\in P$. Thus, $PR' \cap R\subseteq P$.
\end{proof}

\begin{prop}[Incomparability]\label{incomp}
    Let $R\subseteq R'$ be an integral ring extension. If $P'$ and $Q'$ are distinct prime ideals in $R'$ with $P'\cap R= Q' \cap R$ then $P'\not\subseteq Q'$ and $Q'\not\subseteq P'$.
\end{prop}

Geometrically, we are saying that for the map $\Spec R' \to \Spec R$, if points $P',Q'$ are mapped to the same point, then neither is contained in each other---i.e., $P'$ doesn't correspond to an irreducible closed subscheme containing $Q'$ and vice versa. In particular, the fiber of a closed point does not contain a closed subscheme, so we have some sort of finite fiber condition here.

\begin{proof}
    Assume $P' \cap R = Q' \cap R$ and $P'\subseteq Q'$. We prove that $Q'\subseteq P'$ so that $P'=Q'$.
    
    Suppose, to the contrary, that there exists $a\in Q'\setminus P'$. Then, we consider $\overline{a}\in R'/P'$, which is integral over $R/(P'\cap R)$ by lemma \ref{intquot}. Our strategy is to show that this relation on $\overline{a}$ has to be trivial, so $\overline{a}=0$ and $a\in P'$.
    
    Take a relation of minimal degree, and since $a\not\in P'$ the degree is at least 2. We have
    \begin{align*}
        \overline{a}^n + \overline{c_{n-1}}\overline{a}^{n-1} + \cdots + \overline{c_1} \overline{a} + \overline{c_0} &= 0\\
        \overline{a}^n + \overline{c_{n-1}}\overline{a}^{n-1} + \cdots + \overline{c_1} \overline{a} &= -\overline{c_0}
    \end{align*}
    for some $c_0,\dots,c_{n-1} \in R$. Since $a\in Q'$, the left hand side is in $Q'/P'$, so $\overline{c_0}\in Q'/P'$. Since $c_0\in R$, we get
    \[c_0 \in (R \cap Q')/(R\cap P') = 0\]
    (by our original assumption). Thus, we have no constant term in the relation, but since we are in an integral domain $R'/P'$, we can cancel out a factor of $\overline{a}$ to get a monic relation of strictly smaller degree.
    
    Thus, we have a contradiction, so $a\in P'$ and $P'=Q'$.
\end{proof}

\begin{cor}\label{intfield}
    Let $R\subseteq R'$ be an integral ring extension of integral domains. Then, $R$ is a field if and only if $R'$ is a field.
\end{cor}

\begin{proof}
    Assume that $R$ is a field. Let $P'\teq R'$ be a maximal ideal. Then, $P'$ must contract to the $0$ ideal, since the contraction is a prime ideal. Therefore,
    \[0\cap R = P' \cap R,\]
    so by incomparability, $P'=0$. Thus, $R'$ has only maximal ideal $0$, so it is a field.
    
    Now, assume that $R$ is not a field, so there exists a non-zero maximal ideal $\m \teq R$. By Lying Over proposition \ref{lyingover}, there exists a prime ideal $P'\teq R'$ with $P'\cap R = \m$. In particular, $P' \neq 0$, so $R'$ has a nontrivial proper ideal and is not a field.
\end{proof}

\section{Normalization}

\begin{rmk}
    The idea of Noether Normalization is very nice. In one sentence, for any affine variety, we can pick coordinates $z_1,\dots,z_r$ such that projection is a finite morphism. See Gathmann's Example 10.1.
\end{rmk}

\begin{prop}[Noether Normalization] \label{noether}
    Let $R$ be a finitely generated algebra over a field $K$ with generators $x_1,\dots,x_n \in R$. Then there is an injective $K$-algebra homomorphisms $K[z_1,\dots,z_r]\to R$ from a polynomial ring over $K$ to $R$ that makes $R$ into a finite extension ring of $K[z_1,\dots,z_r]$.
    
    Moreover, if $K$ is an infinite field the images of $z_1,\dots,z_r$ in $R$ can be chosen to be $K$-linear combinations of $x_1, \dots, x_n$.
\end{prop}

We first give the proof idea, before going to some necessary lemmas and the actual detailed proof. The idea is as follows: Pick generators $x_1,\dots,x_n$. Check if there is an algebraic relation among $x_1,\dots,x_n$, i.e., a nonzero polynomial $f$ in $x_1,\dots,x_n$ over $K$ with $f(x_1,\dots,x_n)=0$. If there is no algebraic relation, then $R$ is isomorphic the polynomial ring $K[z_1,\dots,z_n]$. Otherwise, we use some lemmas to choose new coordinates $y_1,\dots,y_n$ such that $y_n$ is redundant, in the sense that $y_n$ is integral over $K[y_1,\dots,y_{n-1}]\subseteq R$. We keep adjusting our coordinates until $y_1,\dots,y_r$ have no algebraic relations, in which case we have found our copy of the polynomial ring $K[z_1,\dots,z_r]$ inside $R$.

We give two lemmas to showcase two strategies for picking new coordinates---one that is linear but only works for infinite fields, and one that works in general.

\begin{lem}[Linear Coordinate Change for Infinite Fields] \label{newcoordsinf}
    Let $f\in K[x_1,\dots,x_n]$ be a nonzero polynomial over an infinite field $K$. Then there exists $\lambda\in K$ (to get rid of scalar) and $a_1,\dots,a_{n-1}\in K$ such that
    \[\lambda f(y_1+a_1 y_n, y_2 + a_2 y_n, \dots, y_{n-1} + a_{n-1}y_n, y_n) \quad \in K[y_1,\dots,y_n]\]
    is monic in $y_n$ (i.e., as an element of $(K[y_1,\dots,y_{n-1}])[y_n]$).
\end{lem}

\begin{proof}
    Let $\lambda, a_1,\dots,a_n \in K$ be arbitrary. We compute $f$ written in our new coordinates, and then we determine how to choose our parameters. Writing
    \[f = \sum_{k_1,\dots,k_n} c_{k_1,\dots,k_n} x_1^{k_1}\cdots x_n^{k_n},\]
    we get
    \begin{align*}
        \lambda f(y_1+a_1 &y_n, y_2 + a_2 y_n, \dots, y_{n-1} + a_{n-1}y_n, y_n)\\
        &= \lambda \sum_{k_1,\dots,k_n} c_{k_1,\dots,k_n} (y_1+a_1y_n)^{k_1}\cdots (y_{n-1}+a_{n-1}y_n)^{k_{n-1}} y_n^{k_n}.
    \end{align*}
    Let $d=\deg f$. Then, the maximum $y_n$ degree is $d$. Compute an expression for the $y_n$ degree $d$ part, we get
    \[
        \lambda \sum_{\substack{k_1,\dots,k_n\\ k_1+\cdots+k_n=d}} c_{k_1,\dots,k_n} a_1^{k_1}\cdots a_{n-1}^{k_{n-1}} y_n^{k_1+\dots+k_n} = \lambda f_d(a_1,\dots,a_{n-1},1) y^d,
    \]
    where $f_d$ is the degree $d$ part of $f$. We can therefore guarantee $f$ in the new coordinates to be monic, as long as we can choose $f_d(a_1,\dots,a_{n-1},1) \in K$ nonzero and $\lambda$ to be the inverse. The next lemma guarantees such a choice of $a_1,\dots,a_{n-1}$ is possible.
\end{proof}

\begin{lem}
    Let $f\in K[x_1,\dots,x_n]$ be a nonzero homogeneous polynomial over an infinite field $K$. Then, there are $a_1,\dots,a_{n-1}\in K$ such that $f(a_1,\dots,a_{n-1},1)\neq 0$.
\end{lem}

\begin{proof}
    We prove the lemma by induction on $n$. If $n=1$, then $f=cx$ for $c\in K$ nonzero, so $f(1)=c$. Now, assume $n>1$. We can view $f$ as a polynomial in $(K[x_2,\dots,x_n])[x_1]$ and write
    \[f=\sum_{i=0}^d f_i x_1^i,\]
    where $d$ is the degree of $f$, and the (potentially zero) $f_i$ are homogeneous degree $d-i$ polynomials in $K[x_2,\dots,x_n]$. At least one $f_i$ is nonzero because $f$ is nonzero. Then, by induction we can choose $a_2,\dots,a_{n-1}\in K$ such that for this $i$,
    \[f_i(a_2,\dots,a_{n-1},1)\neq 0.\]
    Fixing these $a_2,\dots,a_{n-1}\in K$, we get
    \[f(x_1,a_2,\dots,a_{n-1},1) \in K[x_1]\]
    a nonzero polynomial. It has at most finitely many roots (i.e., bounded by degree), and since $K$ is infinite, we can choose $x_1=a_1$ to not be a root.
\end{proof}

\begin{lem}[Coordinate Change for Arbitrary Fields] \label{newcoordsfin}
    Let $f\in K[x_1,\dots,x_n]$ be a nonzero polynomial over an arbitrary field $K$. Then there exists $\lambda\in K$ and $a_1,\dots,a_{n-1}\in \N$ such that
    \[\lambda f(y_1+ y_n^{a_1}, y_2 + y_n^{a_2}, \dots, y_{n-1} + y_n^{a_{n-1}}, y_n) \quad \in K[y_1,\dots,y_n]\]
    is monic in $y_n$ (i.e., as an element of $(K[y_1,\dots,y_{n-1}])[y_n]$).
\end{lem}

\begin{proof}
    We prove the claim by induction on $n$. For $n=1$, $f=c_n x^n + \cdots + c_1 x + c_0$. We can then choose $\lambda = \frac{1}{c_n}$.
    
    Now, assume that $\deg f > 1$. Let $d=\deg_{x_1} f$. Then, we can write
    \[f = gx_1^d + h,\]
    where $g\in K[x_2,\dots,x_n]$ and $h\in K[x_1,\dots,x_n]$ with $\deg_{x_1} h \leq d-1$. By our induction hypothesis, we can choose $\lambda \in K$ and $a_2,\dots,a_{n-1}\in \N$ such that
    \[\lambda g(y_2+y_n^{a_2},\dots,y_{n-1}+y_n^{a_n}, y_n)\]
    is monic in $y_n$. Denote its leading $y_n$ term by $y_n^k$. Fix $a_1$, we will see how to choose it later. We write $\tilde{f}$ to denote the evaluation of $f$ at our new coordinates, i.e., the image of $f$ under the homomorphism
    \begin{align*}
        \phi: K[x_1,\dots,x_n] &\longrightarrow K[y_1,\dots,y_n]\\
        x_i &\longmapsto y_i + y_n^{a_i} \qquad i\neq n\\
        x_n &\longmapsto y_n.
    \end{align*}
    Then, we have
    \[\tilde{f} = \tilde{g}(y_1+y_n^{a_1})^d + \tilde{h}.\]
    The leading $y_n$ term of $\tilde{g}(y_1+y_n^{a_1})$ is just $\frac{1}{\lambda} y_n^{k+a_1 d}$. On the other hand, writing the $y_n$ degree of $\tilde{h}$ is bounded a function of $a_2,\dots,a_{n-1}$ plus a term $a_1 (d-1)$. Therefore, we can choose $a_1$ sufficiently large so $\frac{1}{\lambda} y_n^{k+a_1 d}$ dominates $\tilde{h}$ in $y_n$ degree (e.g., choose $a_1> a_2\cdots a_{n-1}\deg f$), giving us a choice of $a_1,\dots,a_{n-1}\in \N$ and $\lambda\in K$.
\end{proof}

Finally, with our two strategies for picking coordinates, we are in a position to prove Noether Normalization.

\begin{proof}[Proof of Noether Normalization \ref{noether}]
    By induction on the given number of generators, $n$. For $n=0$, we get $R=K$ and we are done. 
    
    Otherwise, assume $n>0$. Consider the map
    \begin{align*}
        K[z_1,\dots,z_n] &\longrightarrow R\\
        z_i &\longmapsto x_i
    \end{align*}
    from the polynomial ring to $R$. If the map is injective (i.e., no algebraic relations among the $x_i$), then $R\cong K[z_1,\dots,z_r]$ for $r=n$, and we are done.
    
    Otherwise, the map is not injective and there is a nonzero relation $f\in K[z_1,\dots,z_n]$, i.e., $f(x_1,\dots,x_n)=0$. Using either lemma \ref{newcoordsinf} or lemma \ref{newcoordsfin} (depending on the size of $K$), pick new generators $y_1,\dots,y_n$, i.e.,
    \begin{itemize}
        \item For lemma \ref{newcoordsinf}, choose $y_i \defeq x_i - a_i x_n$ for $i\neq n$, and $y_n \defeq x_n$.
        \item For lemma \ref{newcoordsfin}, choose $y_i \defeq x_i - x_n^{a_i}$ for $i\neq n$, and $y_n \defeq x_n$.
    \end{itemize}
    
    Now, these new generators satisfy the algebraic relation guaranteed by the lemmas, which is monic in $y_n$. This shows that $y_n$ is integral over the subalgebra $K[y_1,\dots,y_{n-1}]$. Therefore, $R=K[y_1,\dots,y_{n-1}][y_n]$ is finite over the subalgebra $K[y_1,\dots,y_{n-1}]$. By induction, there exists $r\in \N$ and an injective homomorphism
    \[K[z_1,\dots,z_r] \longrightarrow K[y_1,\dots,y_{n-1}]\]
    such that $K[y_1,\dots,y_{n-1}]$ is finite over the image. Since finiteness is transitive by proposition \ref{transint}, we get that $R$ is finite over the image. This completes the induction.
    
    If $K$ is infinite, we will always do linear coordinate changes to get the images of $z_1,\dots,z_r$ to be $K$-linear combinations of $x_1,\dots,x_n$.
\end{proof}

\section{Nullstellensatz}

\begin{cor}[Zariski's Lemma]\label{zariski}
    Let $K$ be a field, and let $R$ be a finitely generated $K$-algebra which is also a field.
    Then, $K\subseteq R$ is a finite field extension. In particular, if $K$ is algebraically closed, then $R=K$.
\end{cor}

\begin{proof}
    By Noether Normalization in proposition \ref{noether}, we know $R$ is finite over a polynomial ring $K[z_1,\dots,z_n]$, so it is in particular integral over $K[z_1,\dots,z_n]$ by proposition \ref{integral}. Since $R$ is a field and the extension is integral, by corollary \ref{intfield}, $K[z_1,\dots,z_n]$ must be field. Thus, $r=0$ for $K[z_1,\dots,z_r]$ to be a field (e.g., for $z_1$ to be invertible).
    
    Thus, $K\subseteq R$ is a finite field extension, and if $K$ is algebraically closed, this implies $R=K$.
\end{proof}

\begin{cor}[Weak Nullstellensatz]\label{weak}
    Let $K$ be an algebraically closed field. Then all maximal ideals of the polynomial ring $K[x_1,\dots,x_n]$ are of the form
    \[I(a) = (x_1-a_1, \dots, x_n-a_n)\]
    for some $a=(a_1,\dots,a_n)\in \A_K^n$.
\end{cor}

\begin{proof}
    Let $\m\teq K[x_1,\dots,x_n]$ be a maximal ideal. Then, $K[x_1,\dots,x_n]/\m$ is a field as well as a finitely generated $K$-algebra (use cosets $x_1,\dots,x_n$). Thus, by Zariski's lemma \ref{zariski}, $K[x_1,\dots,x_n]/\m = K$, i.e. identifying $K$ with the image of $K\subseteq K[x_1,\dots,x_n]$ under the quotient map. Then, taking $a_i$ to be the image of $x_i$ under quotienting, we get
    \[(x_1-a_1,\dots,x_n-a_n) \subseteq \m,\]
    and since the left hand ideal is already maximal, we must have
    \[\m = I(a) = (x_1-a_1,\dots,x_n-a_n).\]
\end{proof}

\begin{cor}\label{radpoly}
    For every ideal $I$ in a finitely generated algebra $R$ over a field $K$, we have
    \[\sqrt{I} = \bigcap_{\substack{\m \text{ maximal}\\ \m \supseteq I}} \m.\]
\end{cor}

This is an upgraded version of a theorem about prime ideals instead of maximal ideals. In the language of schemes, we can now figure out if a function is nilpotent by only looking at closed points.

\begin{proof}
    The inclusion $\subseteq$ is immediate because maximal ideals are prime, so they are radical, so $I\subseteq \m$ implies $\sqrt{I}\subseteq \sqrt{\m} =\m$.
    
    Now, we prove the $\supseteq$ inclusion. Let $f\not\in \sqrt{I}$. We want to find $\m \supseteq I$ (a closed point in $V(I)$) such that $f\not\in \m$ ($f$ is nonzero at $\m$). Thus, we want to find a closed point in $D(f)\cap V(I)$. Taking $S=\{1,f,f^2,\cdots\}$, equivalently we are looking for $\m$ such that $\m\cap S =\emptyset$ and $f\not\in \m$.
    
    Since $f\not\in \sqrt{I}$, we know $S\cap I =\emptyset$, so by proposition \ref{getpoint}, we can at least find a prime ideal $P$ with $P\cap S=\emptyset$, $f\not\in P$, and $P_f$ maximal. We now try to show $P$ is in fact maximal anyways.
    
    We consider the field extension which is the composition $K\to R/P \to (R/P)_f = R_f/P_f$. This is indeed a field extension (i.e., not the zero map) because the second map is injective, since $R/P$ is an integral domain.
    
    Then, by construction $R_f/P_f$ is a field, and it is a finitely generated $K$-algebra because it has generators consisting of those of $R\hookrightarrow R_f$ together with $\frac{1}{f}$. Therefore, by Zariski's lemma \ref{zariski}, the extension is finite and in particular integral by proposition \ref{integral}. Thus, $R/P \subseteq R_f/P_f$ is integral as well, so $R/P$ is a field by corollary \ref{intfield}. Therefore, $P$ is maximal.
\end{proof}

\begin{thm}[Hilbert's Nullstellensatz]
    Let $K$ be algebraically closed. Then for every ideal $I\teq K[x_1,\dots,x_n]$ we have $I(V(I)) = \sqrt{I}$.
    
    In particular, there is a bijection
    \begin{align*}
        \{\text{affine varieties in } \A_K^n\} \quad &\longleftrightarrow \quad \{\text{radical ideals in } K[x_1,\dots,x_n]\}\\
        Y \quad &\longmapsto \quad I(Y)\\
        V(I) &\longmapsfrom I.
    \end{align*}
\end{thm}

\begin{proof}
    \textbf{Hard Part.} We prove $I(V(I)) \subseteq \sqrt{I}$. Assume that $f\not\in \sqrt{I}$. We want to find a point $a \in V(I)$ such that $f(a)\neq 0$. By corollary \ref{radpoly}, \[\sqrt{I} = \bigcap_{\substack{\m \text{ maximal}\\ \m \supseteq I}} \m,\]
    so there exists $\m$ maximal such that $f \not\in \m$. By corollary \ref{weak}, this maximal ideal is $I(a)$ for some $a\in K^n$. Since $I\subseteq I(a)$, we have $a\in V(I)$. Since $f\not\in I(a)$, we have $f(a)\neq 0$. Therefore, $f\not\in \sqrt{I}$.
    
    \textbf{Easy Part.} Let $f\in \sqrt{I}$, so $f^n \in I$. Let $a\in V(I)$. Since $f^n\in I$, we know $f^n(a)=0$. Thus, $f(a)=0$, so $f \in I(V(I))$.
    
    This gives us $I(V(I)) = \sqrt{I}$. We also have $V(I(X))=X$ for any affine variety in $\A_K^n$ because clearly $X\subseteq V(I(X))$. Then conversely, we know that since $X$ is an affine variety, $X=V(J)$, so then we want to show $V(J) \supseteq V(I(V(J)))$. By our previous discussion, $J \subseteq \sqrt{J}\subseteq I(V(J))$, and applying $V$ we get our desired subset relation.
    
    Therefore, for radical ideals $I$ and affine varieties $X$, we have $I(V(I)) = I$ and $V(I(X)) = X$. Therefore, $V(-),I(-)$ are inverse bijections.
\end{proof}


\end{document}
